\documentclass[11pt,a4paper]{article}

\usepackage[a-1b]{pdfx}

\usepackage{fontspec}
\usepackage{geometry}
\usepackage{colortbl}
\usepackage{booktabs}  %% \midrule bringt die vertikalen Abstände selbst mit
\usepackage{tabularx}
\usepackage{array}

\setmainfont{texgyreheros}[
  Ligatures      = TeXOff,
  Extension      = .otf,
  UprightFont    = *-regular,
  BoldFont       = *-bold,
  ItalicFont     = *-italic,
  BoldItalicFont = *-bolditalic,
]

\definecolor{akzent}{HTML}{8D313A}
\definecolor{akzentweb}{HTML}{8B1A2B}
\definecolor{akzentalt}{HTML}{8E4848}
\definecolor{haupttext}{HTML}{1C1C1E}
\definecolor{gedaempft}{HTML}{5C5C62}
\definecolor{linie}{HTML}{C9C9CF}

\arrayrulecolor{linie}

\geometry{a4paper, left=25mm, right=20mm, top=18mm, bottom=16mm}

\color{haupttext}
\setlength{\parindent}{0pt}
\pagestyle{empty}

\raggedright

\hyphenpenalty=10000
\exhyphenpenalty=10000
\tolerance=2000

\frenchspacing


\newlength{\balkenhoehe}

\newlength{\positionsabstand}\setlength{\positionsabstand}{2.5mm}
\newlength{\summenabstand}\setlength{\summenabstand}{1mm}
\newlength{\blockabstand}\setlength{\blockabstand}{3mm}
\newlength{\ustidabstand}\setlength{\ustidabstand}{2mm}

\newcommand{\positionstrenner}{\\[\positionsabstand]}
\newcommand{\summentrenner}{\\[\summenabstand]}
\newcommand{\blocktrenner}{\\[\blockabstand]}
\newcommand{\ustidtrenner}{\\[\ustidabstand]}

\newcommand{\zeilenwechsel}{\\}

\newcommand{\kopfperson}[1]{\\{}#1}

\newcommand{\zusammen}[1]{\mbox{#1}}

\newcommand{\rubrik}[1]{{\footnotesize\color{gedaempft}\addfontfeature{LetterSpace=14}\MakeUppercase{#1}}}

\newcommand{\hinweisabsatz}[1]{\par\vspace{6mm}\noindent%
  {\nonfrenchspacing\sfcode`\:=1000 \sfcode`\;=1000 \sfcode`\,=1000 #1}}

\newcommand{\absetzungskopf}[1]{\par\vspace{6mm}\noindent #1\par\vspace{1.5mm}}

\newcommand{\bankblock}[1]{\vspace{9mm}\noindent{\bfseries Bankverbindung}\\[1.5mm]#1}

\newcommand{\absetzungstabelle}[1]{%
  \begin{tabularx}{\linewidth}{@{}>{\raggedright\arraybackslash}X r@{}}#1\end{tabularx}}

\newcommand{\absetzungsdetail}[1]{\newline\hspace*{4mm}{\footnotesize\color{gedaempft}#1}}

\newcommand{\positionsdetail}[1]{\newline{\footnotesize\color{gedaempft} #1}}

\newcommand{\positionszeile}[5]{#1 & #2 & #3 & #4 & #5}

\newcommand{\summenzeile}[2]{#1 & #2}

\newcommand{\abschlusszeile}[2]{\\\midrule\textbf{#1} & \textbf{#2}}

\begin{document}

{\raggedleft\small\color{gedaempft}
  {\bfseries\color{haupttext}Muster IT-Consulting UG (haftungsbeschränkt)}\\
  Beispielweg 42 · 97070 Würzburg\zeilenwechsel{}Tel. +49 931 1234567\zeilenwechsel{}rechnung@muster-it-consulting.example\par}

\vspace{18mm}

\begin{minipage}[t]{0.52\linewidth}\raggedright
  Beispiel Software GmbH\zeilenwechsel{}Industriestraße 1\zeilenwechsel{}90402 Nürnberg
\end{minipage}%
\hfill
\begin{minipage}[t]{0.44\linewidth}
  \raggedleft\small\color{gedaempft}
  Geschäftsführung\zeilenwechsel{}Dr. Erika Musterfrau\blocktrenner{}Leistungszeitraum: \zusammen{01.06.2026\,–\,31.10.2026}\zeilenwechsel{}Zahlbar ohne Abzug bis 19.~November 2026
\end{minipage}

\vspace{14mm}

{\setlength{\fboxsep}{1.4mm}%
\settoheight{\balkenhoehe}{\large\bfseries 0}%
\noindent\colorbox{akzent}{%
  \raisebox{0pt}[\balkenhoehe][0pt]{%
    \makebox[\dimexpr\linewidth-2\fboxsep\relax][s]{%
      \large\bfseries\color{white}Schlussrechnung R26-014\hfill 5.~November 2026}}}}

\vspace{9mm}

\begin{tabularx}{\linewidth}{@{}>{\raggedright\arraybackslash}X r@{\hspace{1.5mm}}l r r@{}}
  \mbox{\rubrik{Beschreibung}} & \mbox{\rubrik{Menge}} & & \mbox{\rubrik{Einzelpreis}} & \mbox{\rubrik{Betrag}} \\[1.5mm]
  \midrule
\positionszeile{Projekt Nord, Gesamtleistung}{1}{pauschal}{20.000,00~€}{20.000,00~€}
\end{tabularx}

\vspace{5mm}

\hfill
\begin{minipage}{0.55\linewidth}
  \begin{tabularx}{\linewidth}{@{}>{\raggedright\arraybackslash}X r@{}}
  \summenzeile{Nettobetrag}{20.000,00~€} \summentrenner{}
  \summenzeile{zzgl. 19~\% USt.}{3.800,00~€}
  \abschlusszeile{Gesamtbetrag}{23.800,00~€}
  \end{tabularx}
\end{minipage}



\vspace{6mm}



\absetzungskopf{Abzüglich der berechneten und vereinnahmten Abschläge:}
\absetzungstabelle{
  \summenzeile{Abschlagsrechnung R26-009 vom 15.06.2026, vereinnahmt am 10.06.2026\absetzungsdetail{netto 4.000,00~€ zzgl. 19~\% USt. 760,00~€}}{−4.760,00~€} \positionstrenner{}
  \summenzeile{Abschlagsrechnung R26-011 vom 20.08.2026, vereinnahmt am 02.09.2026\absetzungsdetail{netto 4.000,00~€ zzgl. 19~\% USt. 760,00~€}}{−4.760,00~€}
  \abschlusszeile{Restbetrag}{14.280,00~€}
}

\vspace{5mm}

Ich bedanke mich für das Vertrauen und die angenehme Zusammenarbeit.

\vspace{5mm}

Mit freundlichen Grüßen\\[7mm]
Dr. Erika Musterfrau

\bankblock{Kontoinhaber: Muster IT-Consulting UG (haftungsbeschränkt)\zeilenwechsel{}Institut: Musterbank Würzburg\zeilenwechsel{}IBAN: DE89 3704 0044 0532 0130 00\zeilenwechsel{}BIC: MUSTDEFFXXX}

\vfill

{\centering\scriptsize\color{gedaempft}
Muster IT-Consulting UG (haftungsbeschränkt) · Beispielweg 42, 97070 Würzburg · Registergericht Würzburg, HRB 12345\\
USt-IdNr. DE 123 456 788 · Musterbank Würzburg · BIC MUSTDEFFXXX · IBAN DE89 3704 0044 0532 0130 00\par}

\end{document}
